\section{Mechanism: Close Substitutes}

Our central hypothesis is that DFL is most useful when prediction errors change
the objective value of the selected solution, not merely its identity. In
packing-style graph problems, feasible solutions can be composed of several
components. If each component has close substitutes, a prediction error may swap
one feasible solution for another while changing the true objective only
slightly.

This mechanism predicts two empirical patterns. First, solution identity can
change even when normalized regret remains small. Second, all-edge prediction
error should be a weak explanation of decision gap compared with error on
decision-changing edges, such as edges in the symmetric difference between the
oracle and selected solutions.

This framing avoids a stronger claim that DFL is unnecessary for all packing
problems. The claim is conditional: when the feasible landscape contains many
near-optimal alternatives, DFL has less room to improve over a strong two-stage
model; when prediction errors redirect a coupled decision to a much worse
solution, DFL can matter substantially.
